\section{A Quick Tour through Networking}
\subsection{Distributed Systems (DS)}
\subsubsection*{Definition}
\fatsf{Ein verteiltes System ist ein Zusammenschluss mehrere Computer,} die über ein Netzwerk miteinander verbunden sind und \fatsf{als ein einziges System} erscheinen und arbeiten
\subsubsection*{Zitat zur Veranschaulichung}
\fatsf{Laslie Lamport:} \enquote{Ein verteiltes System ist eines, bei dem der Ausfall eines Computers, von dessen Existenz man nicht einmal wusste, den eigenen Computer unbenutzbar machen kann}\\
Das zeigt: \fatsf{Ein Fehler irgendwo kann das ganze System betreffen}
\subsubsection*{Eigenschaften eines verteilten Systems}
\begin{itemize}
    \item Besteht aus \fatsf{mehreren Computern} (oft heterogen)
    \item Verbindet diese über ein \fatsf{Computernetzwerk}
    \item Die Computer \fatsf{kooperieren eng}, um \fatsf{gemeinsam eine Aufgabe oder einen Dienst zu erfüllen}
    \item Das System funktioniert \fatsf{nur über das Netzwerk}
    \item Das System wirkt \fatsf{nach außen wie ein einziges} (\enquote{System of systems})
\end{itemize}
\subsubsection*{Herausforderungen/Probleme}
\begin{itemize}
    \item Verzögerungen im Netzwerk (Latenz)
    \item Ausfall einzelner Komponenten
    \item Datenkonsistenz und Synchronisation
    \item Fehlende zentrale Kontrolle
\end{itemize}
\subsubsection*{Beispiel für verteilte Systeme}
\begin{itemize}
    \item \fatsf{Internet}
    \item \fatsf{IPTV} (Fernsehen über Internet)
    \item \fatsf{VoIP} (z.B. Skype, Zoom)
    \item \fatsf{Peer-to-Peer Netzwerke} (z.B. BitTorrent)
    \item \fatsf{Online-Spiele}
    \item \fatsf{Metaverse / Virtuelle Welten}
\end{itemize} 
\subsubsection*{Merksatz}
\enquote{Kein verteiltes System ohne Computernetzwerk}
\subsection{Netzwerke im Überblick}
\subsubsection*{Allgemeines}
\begin{itemize}
    \item Netzwerke sind ein \fatsf{junges, schnell wachsendes Fachgebiet}
    \item Gleichzeitig aber auch alltäglich geworden: eine \fatsf{\enquote{Commodity} bzw. \enquote{Utility}} (wie Strom oder Wasser)
\end{itemize}
\subsubsection*{Ständige Veränderung (Constant Flux)}
\begin{itemize}
    \item \fatsf{Rasante technologische Entwicklung:}
    \begin{itemize}
        \item \fatsf{Kehrnetzwerke:} Von T1 (1,544 Mbps) $\rightarrow$ zu 100-Gbps-Switches (und mehr)
        \item \fatsf{Heimanwender:} 1,2 kbps Modem $\rightarrow$ ADSL2 + (24 Mbps) $\rightarrow$ VDSL2-VPlus (300 Mbps) $\rightarrow$ Glasfaser (n x 1 Gbps)
    \end{itemize}
    \item \fatsf{Wandel in der Infrastruktur und Kommunikation:}
    \begin{itemize}
        \item Geräte: Workstations $\rightarrow$ PCs $\rightarrow$ Smartphones $\rightarrow$ Sensoren $\rightarrow$ Kommunikation \& Messung kombiniert
        \item Betriebssysteme: Unix $\rightarrow$ Linux/BSD $\rightarrow$ Windows, macOS $\rightarrow$ iOS, Android, TinyOS $\rightarrow$ \fatsf{Jedes Gerät hat heute ein OS}
    \end{itemize}
\end{itemize}
\subsubsection*{Vergleich: Programmiersprachen vs. Netzwerktechnik}
\begin{table}[H]
    \centering
    \rowcolors{2}{\thepagecolor}{fgcolor!10!\thepagecolor}
    \renewcommand{\arraystretch}{1.2}
    \begin{tabular}{lll}
        \toprule
        \fatsf{Bereich} & \fatsf{Programmiersprachen} & \fatsf{Netzwerke / Verteilte Systeme (CN \& DS)} \\
        \midrule
        Theorie & Klar definierte Grundlage & Viele verschiedene Modelle\\
        Struktur & Gut organisierte Forschung & Forschung eher unstrukturiert\\
        Reife & Mehrere Phasen der Reifung durchlaufen & Teilweise noch in den Kinderschuhen?\\
        Realität & Softwarequalität oft sehr schlecht & Große funktionierende Systeme existieren\\
        \bottomrule
    \end{tabular}
\end{table}
\subsubsection*{Merksätze}
\begin{itemize}
    \item \enquote{Networking ist jung, aber überall}
    \item \enquote{Netze verändern sich ständig - Geräte, Dienste, Kommunikation, Betriebssysteme}
\end{itemize}
\clearpage
\subsection{Kommunikation zwischen zwei Geräte}
\subsubsection*{Beispiel 1: World Wide Web}
Was passiert, wenn du \texttt{www.tu-darmstadt.de} in die Adresszeile eingibst und Enter drückst?
\begin{enumerate}
    \item \fatsf{Namensauflösung (DNS):} Domainname $\rightarrow$ IP-Adresse
    \item \fatsf{Verbindungsaufbau:} z.B. über TCP
    \item \fatsf{Anfrage (Request):} Browser sendet HTTP-Request
    \item \fatsf{Antwort (Response):} Webserver liefert die Webseite zurück
    \item \fatsf{Darstellung:} Browser zeigt die Seite an
\end{enumerate}
$\rightarrow$ \fatsf{Übertragung einer Datei oder von Datenpaketen}
\subsubsection*{Beispiel 2: Telefonanruf}
\fatsf{Was passiert, wenn du jemanden anrufst?}
\begin{enumerate}
    \item \fatsf{Teilnehmernummer finden \& Verbindungsaufbau}
    \item \fatsf{Sprachdaten in Echtzeit übertragen} (z.B. über VoIP oder klassische Telefonleitung)
    \item \fatsf{Beidseitiger, synchroner Austausch}
\end{enumerate}
$\rightarrow$ \fatsf{Kontinuierlicher Informationsfluss - muss rechtzeitig ankommen!}
\subsubsection*{Vergleich Webseite vs. Telefonat}
\begin{table}[H]
    \centering
    \rowcolors{2}{\thepagecolor}{fgcolor!10!\thepagecolor}
    \renewcommand{\arraystretch}{1.2}
    \begin{tabular}{lll}
        \toprule
        \fatsf{Merkmal} & \fatsf{Webseite (HTTP)} & \fatsf{Telefonat (VoIP / klassisch)}\\
        \midrule
        Datenart & \fatsf{Statische Datenmenge} (z.B. HTML-Datei) & \fatsf{Kontinuierlicher Strom} (Sprache)\\
        Zeitkritisch? & Nein (Daten können gepuffert werden) & Ja (muss in \fatsf{Echtzeit} ankommen)\\
        Interaktivität & Einseitig oder zeitverzögert & Echtzeit, \fatsf{bidirektional}\\
        \bottomrule
    \end{tabular}
\end{table}
\subsubsection*{Was ist eigentlich \enquote{Daten} und \enquote{Information}?}
\begin{table}[H]
    \centering
    \rowcolors{2}{\thepagecolor}{fgcolor!10!\thepagecolor}
    \renewcommand{\arraystretch}{1.2}
    \begin{tabular}{ll}
        \toprule
        \fatsf{Begriff} & \fatsf{Erklärung}\\
        \midrule
        \fatsf{Daten} & Rohwerte, z.B. Zahlen, Zeichen - ohne Kontext noch bedeutungslos\\
        \fatsf{Information} & Wenn Daten in einem Kontext \fatsf{interpretiert und verstanden} werden\\
        \bottomrule
    \end{tabular}
\end{table}
Beispiel:
\begin{itemize}
    \item \enquote{42} = \fatsf{Daten}
    \item \enquote{42 Grad Celsius Körpertemperatur} = \fatsf{Information}
\end{itemize}
\subsubsection*{Einfachste Kommunikation: Direkte physikalische Verbindung}
\fatsf{Beispiele:}
\begin{itemize}
    \item \fatsf{WWW:} Client (Browser) direkt per Kabel mit Server verbunden
    \begin{itemize}
        \item Client sendet Anfrage, Server liefert Daten
    \end{itemize}
    \item \fatsf{Telefonie:} Zwei Telefone direkt per Kabel verbunden
\end{itemize}
\subsubsection*{Wichtige Begriffe: Information - Daten - Signale}
\begin{table}[H]
    \centering
    \rowcolors{2}{\thepagecolor}{fgcolor!10!\thepagecolor}
    \renewcommand{\arraystretch}{1.2}
    \begin{tabular}{ll}
        \toprule
        \fatsf{Abstrakte Ebene} & \fatsf{Reale Umsetzung}\\
        \midrule
        \fatsf{Information} & Bedeutung, Inhalt\\
        \fatsf{Daten} & Darstellung der Information (z.B. als Bits)\\
        \fatsf{Signale} & Physikalische Repräsentation von Daten (z.B. elektrische Spannung, Schall)\\
        \bottomrule 
    \end{tabular}
\end{table}
\subsubsection*{Was ist ein Signal?}
Ein \fatsf{Signal} ist die Darstellung von Daten durch charakteristische Änderungen eine physikalischen Werts (über Raum oder Zeit)
\begin{itemize}
    \item \fatsf{Materiell:} Buchstaben auf Papier
    \item \fatsf{Immateriell:} Schallwellen beim Sprechen
    \item \fatsf{Elektrisch:} Spannung oder Strom im Draht
\end{itemize}
\subsubsection*{Bits und Signale}
\begin{table}[H]
    \centering
    \rowcolors{2}{\thepagecolor}{fgcolor!10!\thepagecolor}
    \renewcommand{\arraystretch}{1.2}
    \begin{tabular}{ll}
        \toprule
        Was soll übertragen werden? & Daten (als Bits)\\
        Was kann physisch übertragen werden? & Signale (z.B. Spannung)\\
        \bottomrule 
    \end{tabular}
\end{table}
\fatsf{Wir brauchen:}
\begin{itemize}
    \item \fatsf{Kodierung:} Bits $\rightarrow$ Signale (Sender)
    \item \fatsf{Dekodierung:} Signale $\rightarrow$ Bits (Empfänger)
\end{itemize}
\fatsf{Einfaches Beispiel (nicht praxisnah):}
\begin{itemize}
    \item Bit 1 = Strom vorhanden
    \item Bit 0 = Kein Strom
\end{itemize}
\subsubsection*{Mekrsätze}
\begin{itemize}
    \item \enquote{Daten werden zu Signalen, damit Maschinen miteinander sprechen können}
    \item \enquote{Ein Kabel speichert Daten, indem es sie unterwegs hält}
    \item \enquote{Nur mit Signalen erreichen Bits ihr Ziel}
\end{itemize}
\subsubsection*{Technische Eigenschaften der Übertragung}
\begin{table}[H]
    \centering
    \rowcolors{2}{\thepagecolor}{fgcolor!10!\thepagecolor}
    \renewcommand{\arraystretch}{1.2}
    \begin{tabular}{ll}
        \toprule
        \fatsf{Begriff} & \fatsf{Erklärung}\\
        \midrule
        \fatsf{Verzögerung (Delay d)} & Wie lange braucht ein Signal\\
        \fatsf{Signalausbreitung (v)} & Geschwindigkeit der Signalverbreitung (z.B. Lichtgeschwindigkeit \texttt{c}, in Kupfer: ca. \texttt{2/3 c})\\
        \fatsf{Datenrate (r)} & Wie viele Bits pro Sekunde können gesendet werden?\\
        \fatsf{Fehlerrate} & Wie viele Bits kommen fehlerhaft an? Welche Fehlerbilder treten auf?\\
        \bottomrule 
    \end{tabular}
\end{table}
\subsubsection*{Wie kann ein Kabel Daten \enquote{speichern}?}
\begin{itemize}
    \item Während der ersten \texttt{d} Sekunden ist das erste Bit \fatsf{unterwegs}, aber noch \fatsf{nicht angekommen}
    \item In dieser Zeit hat der Sender bereits \texttt{d*r} Bits gesendet - diese Bits \fatsf{befinden sich im Kabel $\rightarrow$ \enquote{gespeichert} in der Leitung}
\end{itemize}
\begin{figure}[h]
    \centering
    \includegraphics[width=0.3\linewidth]{Chapter 1_Images/MSC.jpeg}
\end{figure}
\subsubsection*{Zwei-Wege-Kommunikation: Übertragungsarten}
\begin{table}[H]
    \centering
    \rowcolors{2}{\thepagecolor}{fgcolor!10!\thepagecolor}
    \renewcommand{\arraystretch}{1.2}
    \begin{tabular}{lll}
        \toprule
        \fatsf{Art} & \fatsf{Beschreibung} & \fatsf{Beispiel}\\
        \midrule
        \fatsf{Simplex} & Nur eine Seite sendet & Rodio\\
        \fatsf{Halbduplex} & Sender und Empfänger wechseln sich ab & Walkie-Talkie, Gespräch\\
        \fatsf{Vollduplex} & Beide Seiten können gleichzeitig senden & Telefonat, Webkommunikation\\
        \bottomrule 
    \end{tabular}
\end{table}
\clearpage
\subsubsection*{Halbduplex: Wechselnde Kommunikation}
\begin{itemize}
    \item \fatsf{Problem:} Beide senden gleichzeitig $\rightarrow$ Signale stören sich
    \item \fatsf{Frage:} Wie erkennt ein Teilnehmer, wann der andere fertig ist?
\end{itemize}
\fatsf{Lösungen:}
\begin{figure}[h]
    \centering
    \includegraphics[width=0.8\linewidth]{Chapter 1_Images/Half duplex.jpeg}
\end{figure}
\begin{enumerate}
    \item \fatsf{Time Division Duplex (TDD)}\\
    $\rightarrow$ Feste Zeitslots, synchronisiert\\
    Beispiel: A darf 10 ms senden, dann B für 10 ms usw.
    \item \fatsf{On-Demand / dynamisch:}
    \begin{itemize}
        \item Sender kündigt Sendedauer vorher an
        \item Oder: schickt spezielles \enquote{Ende}-Signal
    \end{itemize}
\end{enumerate}
\fatsf{Herausforderungen:}
\begin{itemize}
    \item Max. Länge muss definiert sein
    \item \enquote{Ende}-Signal darf \fatsf{nicht Teil der Nutzerdaten} sein
    \item \fatsf{Bit-Synchronisation} nötig
\end{itemize}
\subsubsection*{Vollduplex - Gleichzeitiges Senden \& Empfangen}
\begin{itemize}
    \item \fatsf{Kabelgebunden:}
    \begin{itemize}
        \item Zwei separate Kabel (je Richtung) $\rightarrow$ zuverlässig, aber aufwändig (Kosten, Platz, Wartung)
        \item \fatsf{Ein Kabel?} Möglich auf kurze Distanz (z.B. mit Signalverarbeitung, z.B. Echo-Unterdrückung)
    \end{itemize}
    \item \fatsf{Drahtlos / Funk / Kabel-TV:}
    \begin{itemize}
        \item \fatsf{Frequenz Division Duplex (FDD):}\\
        Unterschiedliche Frequenzen für Hin- und Rückrichtung\\
        \fatsf{$\rightarrow$ keine Kollision} durch Frequenztrennung
    \end{itemize}
\end{itemize}
\begin{figure}[h]
    \centering
    \includegraphics[width=0.8\linewidth]{Chapter 1_Images/Full duplex.jpeg}
\end{figure}
\subsubsection*{Voll-Duplex mit TDD - Geht das überhaupt?}
\fatsf{Klingt widersprüchlich, aber:}
\begin{itemize}
    \item Wenn \fatsf{Übertragungsmedium mindestens doppelt so schnell} ist wie die benötige Datenrate, dann:
    \begin{itemize}
        \item Teilnehmer A speichert Daten zwischen (z.B. als Datei oder in einem Speicherchip), während B sendet
        \item Danach sendet A \fatsf{zwischengespeicherte + neue} Daten in kürzerer Zeit\\
        $\rightarrow$ So wirkt es wie \fatsf{gleichzeitiges Senden}
    \end{itemize}
\end{itemize}
\begin{figure}[h]
    \centering
    \includegraphics[width=0.8\linewidth]{Chapter 1_Images/TDD.jpeg}
\end{figure}
\begin{itemize}
    \item \fatsf{Funktioniert bei digitalen Signalen}
    \item \fatsf{Nicht geeignet für analoge Signale} (z.B. klassische analoge Telefonie)\\
    $\rightarrow$ Da keine exakte Zwischenspeicherung möglich ist
\end{itemize}
\subsection{Netzwerkstruktur \& Vermittlungsverfahren}
\fatsf{Warum nicht jedes Gerät direkt verbinden?}
\begin{itemize}
    \item Direkte Verbindung \fatsf{jedes Geräts mit jedem anderen} ist \fatsf{nicht skalierbar}\\
    $\rightarrow$ Bei n Geräten: \fatsf{n(n-1)/2} Verbindungen nötig (explodiert schnell!)
    \begin{figure}[h]
        \centering
        \includegraphics[width=0.8\linewidth]{Chapter 1_Images/Beirut Connections.jpeg}
    \end{figure}
    \clearpage
    \item Lösung: \fatsf{Struktur im Netzwerk} einführen
    \begin{figure}[h]
        \centering
        \includegraphics[width=0.55\linewidth]{Chapter 1_Images/Structure into a Network.jpeg}
    \end{figure}
\end{itemize}
\subsubsection*{Netzwerkstruktur}
Zwei Arten von Komponenten:
\begin{itemize}
    \item \fatsf{Endsysteme:} Benutzergeräte (PCs, Smartphones, etc.)
    \item \fatsf{Vermittlungselemente:} Router, Switches
\end{itemize}
\subsubsection*{Vermittlungsarten (Switching)}
\begin{enumerate}
    \item \fatsf{Leitungsvermittlung (Circuit Switching)}
    \begin{itemize}
        \item Eine \fatsf{dedizierte Verbindung} (Schaltkreis) wird \fatsf{vorab} zwischen zwei Endsystemen aufgebaut\\
        $\rightarrow$ Wie klassisches Telefonnetz (analog: \enquote{Fräulein vom Amt})
        \item \fatsf{Verbindung bleibt exklusiv} für die Dauer der Kommunikation bestehen
    \end{itemize}
    \begin{itemize}
        \item \fatsf{Vorteile:}
        \begin{itemize}
            \item Feste Ressourcen $\rightarrow$ \fatsf{garantierte Übertragung}
            \item \fatsf{Keine Verzögerung durch Warteschlangen} (kein Queuing)
        \end{itemize}
        \item \fatsf{Nachteile:}
        \begin{itemize}
            \item Ressourcen \fatsf{blockiert}, auch bei Pausen
            \item Verbindung muss \fatsf{vor dem Datentransfer aufgebaut} werden
            \item Wenig effizient bei kurzen Nachrichten oder vielen Teilnehmern
        \end{itemize}
    \end{itemize}
    \begin{figure}[h]
        \centering
        \includegraphics[width=0.72\linewidth]{Chapter 1_Images/Circuit Switching.jpeg}
    \end{figure}
    \clearpage
    \item \fatsf{Paketvermittlung (Packet Switching)}
    \begin{itemize}
        \item Daten werden in \fatsf{Pakete} aufgeteilt (jeweils mit Zieladresse)
        \item Pakete werden \fatsf{einzeln und unabhängig} durch das Netz geschickt
        \item Kein fester Pfad nötig $\rightarrow$ \fatsf{dynamisch}es Routing möglich
    \end{itemize}
    \fatsf{Eigenschaften:}
    \parpic(6cm,3.5cm)[r]{
        \includegraphics[width=4.5cm, height=6cm]{Chapter 1_Images/store-and-forward.jpeg}
    }
    \begin{itemize}
        \item Jeder \fatsf{Switch} muss:
        \begin{enumerate}
            \item Paket empfangen
            \item Zwischenspeichern (Buffering)
            \item Zieladresse auslesen
            \item Nächsten Weiterleitungspfad bestimmen
            \item Paket weiterleiten (\enquote{store-and-forward})
        \end{enumerate}
    \end{itemize}
    \fatsf{Vorteile:}
    \begin{itemize}
        \item \fatsf{Effiziente Ressourcennutzung}, da keine exklusive Leitung
        \item Flexibel \& robust bei Ausfällen oder Störungen
    \end{itemize}
    \fatsf{Herausforderungen:}
    \begin{itemize}
        \item Start/End eines Pakets erkennen
        \item Adressing und Routing
        \item Pufferüberläufe \& Status (Congestion)
        \item Höherer \fatsf{Verwaltungsaufwand} durch Header \& Steuerinformationen
    \end{itemize}
\end{enumerate}
\subsubsection*{Merksätze}
\begin{itemize}
    \item \enquote{Direkte Verbindungen skalieren nicht - Vermittlung ist nötig}
    \item \enquote{Leitungsvermittlung = reserviert \& einfach, aber unflexibel}
    \item \enquote{Paketvermittlung = flexibel \& effizient, aber komplexer}
\end{itemize}
\clearpage
\subsection{Multiplexing \& Shared Resources}
\fatsf{Multiplexing Definition:}
\begin{itemize}
    \item Organisation der Übertragung mehrere Datenströme über eine einzige physische Verbindung\\
    $\rightarrow$ Verschiedene Sender teilen sich ein Übertragungsmedium
    \item Das Ziel ist es, \fatsf{eine Leitung effizient zu nutzen}, statt für jede Verbindung eine eigene Leitung bereitzustellen
\end{itemize}
\begin{figure}[h]
    \centering
    \includegraphics[width=0.4\linewidth]{Chapter 1_Images/Multiplexing.jpeg}
\end{figure}
\subsubsection*{Wie funktioniert Multiplexing konkret?}
\fatsf{Beispiel: Zwei PCs senden gleichzeitig an denselben Empfänger}
\begin{itemize}
    \item Wenn beide gleichzeitig auf derselben Leitung senden, \fatsf{würden sich ihre Signale überlappen $\rightarrow$ Signal-Kollision!}
    \item Multiplexing sorgt dafür, dass diese Übertragungen \fatsf{geordnet} und \fatsf{kollisionsfrei} ablaufen - entweder in der Zeit, über verschiedene Frequenzen, oder durch andere Techniken
\end{itemize}
\subsubsection*{Multiplexing-Methoden}
\begin{table}[H]
    \centering
    \rowcolors{2}{\thepagecolor}{fgcolor!10!\thepagecolor}
    \renewcommand{\arraystretch}{1.2}
        \begin{tabular}{p{4cm} p{6.5cm} p{4cm}}
        \toprule
        \fatsf{Methode} & \fatsf{Beschreibung} & \fatsf{Beispiel}\\
        \midrule
        \fatsf{TDM} (Time Division Multiplexing) & Jede Verbindung bekommt eine feste \fatsf{Zeitspanne} (Zeitschlitz) zur Übertragung & Telefonnetz, Ethernet \\
        \fatsf{FDM} (Frequency Division Multiplexing) & Daten werden gleichzeitig, aber auf \fatsf{unterschiedlichen Frequenzen} gesendet & UKW-Radio, Kabelfernsehen\\
        \fatsf{WDM} (Wavelength Division Multiplexing) & Wie FDM, aber in der \fatsf{optischen Kommunikation (Glasfaser)} & Glasfasernetze\\
        \fatsf{CDM} (Code Division Multiplexing) & Jeder Sender bekommt einen eindeutigen Code & Mobilfunk (3G/UMTS)\\
        \fatsf{SDM} (Space Division Multiplexing) & Verschiedene Signale laufen über \fatsf{räumlich getrennte Kanäle, z.B. in Mobilfunkzellen} & Mobilfunknetze (4G/5G) \\
        \bottomrule
    \end{tabular}
\end{table}
\subsubsection*{Demultiplexing}
\fatsf{Gegenstück zu Multiplexing:} Trennt am Empfänger die einzelnen Datenströme wieder auf
\subsubsection*{Upward vs. Downward Multiplexing} 
\begin{table}[H]
    \centering
    \rowcolors{2}{\thepagecolor}{fgcolor!10!\thepagecolor}
    \renewcommand{\arraystretch}{1.2}
        \begin{tabular}{l p{6cm} p{6cm}}
        \toprule
        \fatsf{Upward} & Viele logische Verbindungen teilen sich \fatsf{eine physische Verbindung} & Mehrere Apps auf deinem Handy nutzen gleichzeitig das WLAN\\
        \fatsf{Downward} & Eine logische Verbindung nutzt \fatsf{mehrere physische Verbindungen} & Download-Boosting durch mehrere Leitungen\\
        \bottomrule
    \end{tabular}
\end{table}
\subsubsection*{Shared Resources}
\begin{itemize}
    \item Multiplexing regelt den \fatsf{Zugriff auf gemeinsam genutzte Ressourcen} - z.B. ein drahtloses Medium
    \item \fatsf{Beispiele für geteilte Ressourcen:}
    \begin{itemize}
        \item Funkfrequenzen (Luft als Übertragungsmedium)
        \item Gemeinsamer Draht (Kupferkabel)
        \item Ein Netzwerk-Switch oder Routerport
    \end{itemize}
\end{itemize}
\subsubsection*{Das Problem der Mehrfachzugriffe (Multiple Access)}
\begin{itemize}
    \item \fatsf{Broadcast-Medium} (z.B. Luft, Funk, geteiltes Kabel) $\rightarrow$ Daten erreichen alle Teilnehmer
    \item \fatsf{Nur ein Sender zur gleichen Zeit erlaubt} $\rightarrow$ sonst: Kollisionen
    \item Lösung: \fatsf{Medium Access Control (MAC)}
\end{itemize}
\subsubsection*{Zugriffsverfahren (Medium Access Control / MAC)}
\begin{table}[H]
    \centering
    \rowcolors{2}{\thepagecolor}{fgcolor!10!\thepagecolor}
    \renewcommand{\arraystretch}{1.2}
    \begin{tabular}{l l p{8cm}}
        \toprule
        \fatsf{Verfahren} & \fatsf{Beispiel} & \fatsf{Erklärung}\\
        \midrule
        \fatsf{Zentralisiert (Lehrer-Schüler-Prinzip)} & Unterricht & Nur der Lehrer gibt Schülern das Wort (Zugang zum Medium)\\
        \fatsf{Dezentralisiert (Party-Prinzip)} & Gesprächsrunde & Jeder versucht zu sprechen, wartet auf Pause - Kollisionen möglich, dann Rückzug \& Wiederholung\\
        \bottomrule 
    \end{tabular}
\end{table}
\subsubsection*{Forwarding und Next Hop Selection}
\fatsf{Was ist Forwarding?}
\begin{itemize}
    \item Forwarding ist der \fatsf{Prozess, bei dem ein Router entscheidet, wohin ein Datenpaket als Nächstes geschickt werden soll}, damit es sein Ziel erreicht
    \item Es geht \fatsf{nicht} darum, den ganzen Weg zu planen, sondern nur den \fatsf{nächsten Schritt} - den sogenannten \fatsf{Next Hop}
\end{itemize}
\clearpage
\subsection{Netzwerkorganisation \& Abstraktion}
\fatsf{Ziel}: Wie können viele, verschiedene (\enquote{alien}) Computer effizient kommunizieren?
\subsubsection*{Wichtigstes Konzept: Abstraktion}
\begin{itemize}
    \item \fatsf{Abstraktion} vereinfacht komplexe Systeme, indem sie Details verbirgt
    \item Beispiele:
    \begin{itemize}
        \item In der Programmierung von Maschinencode zu High-Level-Sprachen
        \item In Netzwerken: durch \fatsf{Schichtenmodelle} (Layer Models)
    \end{itemize}
\end{itemize}
\subsubsection*{Netzwerk-Schichtenmodelle (Layer Models)}
\fatsf{OSI-Modell (ISO/OSI)}
\begin{itemize}
    \item \fatsf{7 Schichten} (standardisiert, konzeptionell)
    \item Unterteilt Aufgaben der Datenkommunikation in \fatsf{klar definierte Ebenen}
    \item Ziel: \fatsf{Interoperabilität} und \fatsf{Modularität}
\end{itemize}
\fatsf{Internetmodell}
\begin{itemize}
    \item \fatsf{5-Schichten-Modell} (praktischer, realitätsnäher)
    \item Nutzt im Kern die unteren 5 OSI-Schichten
\end{itemize}
\subsubsection*{Kommunikation im Netzwerk: Was ist ein Protokoll?}
Ein \fatsf{Protokoll} definiert \fatsf{Format, Reihenfolge und Regeln} der Nachrichtenübertragung zwischen zwei oder mehreren Geräten\\
\fatsf{Funktionen eines Protokolls:}
\begin{itemize}
    \item Format \& Struktur der Daten (Header, Payload, etc.)
    \item Verhalten beim Senden und Empfangen
    \item Beispiele für menschliche Protokolle: Begrüßung, Telefonanruf (Anruf $\rightarrow$ Antwort $\rightarrow$ Gespräch)
\end{itemize}
\subsubsection*{Schichtenarchitektur (Layered Architecture)}
\parpic(5.5cm,5cm)[r]{
    \includegraphics[width=8cm, height=4cm]{Chapter 1_Images/Layered Architecture.jpeg}
}
\begin{itemize}
    \item Jede Schicht:
    \begin{itemize}
        \item bietet \fatsf{einen Dienst} (Service) für die darüberliegende Schicht
        \item nutzt den Dienst der darunterliegenden Schicht
    \end{itemize}
    \item Kommunikation erfolgt \fatsf{zwischen gleichrangigen Entitäten (Peers)} \\ derselben Schicht
    \item Vorteile:
    \begin{itemize}
        \item Austauschbarkeit von Schichten
        \item Klarheit \& Wartbarkeit
    \end{itemize}
    \item Nachteil: zusätzlicher Overhead (mehr Protokollinformationen)
\end{itemize}
\subsubsection*{Datenbegriffe in den Schichten (Data Units)}
\begin{table}[H]
    \centering
    \rowcolors{2}{\thepagecolor}{fgcolor!10!\thepagecolor}
    \renewcommand{\arraystretch}{1.2}
    \begin{tabular}{ll}
        \toprule
        \fatsf{Begriff} & \fatsf{Bedeutung}\\
        \midrule
        \fatsf{PDU} & Protocol Data Unit - gesamte Einheit, die auf einer Schicht verarbeitet wird\\
        \fatsf{SDU} & Service Data Unit - \enquote{Nutzerdaten}, die von einer höheren Schicht übergeben wurden\\
        \fatsf{PCI} & Protocol Control Information - Steuerdaten für die Kommunikation auf derselben Ebene\\
        \fatsf{Packet} & Allgemeiner Begriff für Transporteinheit\\
        \fatsf{Datagram} & Verbindungslos gesendetes Packet\\
        \fatsf{Frame} & Datenpaket auf Sicherungsschicht mit finalem Header/Trailer\\
        \fatsf{Cell} & Feste kleine Pakete (z.B. bei ATM)\\
        \bottomrule 
    \end{tabular}
\end{table}
\subsubsection*{Verteilte Systeme \& Netzwerke}
\fatsf{Verteilte Maschine:}
\begin{itemize}
    \item Besteht aus \fatsf{mehreren autonomen Systemen (AS)}, verbunden durch ein \fatsf{Kommunikationssystem (CSS)}
\end{itemize}
\begin{figure}[h]
    \centering
    \includegraphics[width=0.3\linewidth]{Chapter 1_Images/Distributed Systems.jpeg}
\end{figure}
\fatsf{Wichtige Eigenschaft:}
\begin{itemize}
    \item Das Kommunikationssystem \fatsf{versteckt Komplexität} $\rightarrow$ Fehler, Heterogenität, Standorte, etc. werden \fatsf{transparent}
    \item Verteilte Systeme setzen auf \fatsf{funktionierenden Netzwerken auf}
\end{itemize}
\clearpage
\subsection{OSI-Modell}
\fatsf{Was ist das OSI-Modell?}
\begin{itemize}
    \item Das \fatsf{OSI-Modell} ist ein \fatsf{theoretisches Modell}, das die Kommunikation zwischen zwei Computern in \fatsf{sieben Schichten (Layers)} aufteilt
    \item Jede Schicht hat \fatsf{eine bestimmte Aufgabe} und arbeitet mit der darunterliegenden zusammen - wie bei einer \fatsf{Leiter}, von der Hardware bis zur Software
\end{itemize}
\subsubsection*{Warum OSI?}
\begin{itemize}
    \item \fatsf{Vereinheitlichung:} Damit Geräte unterschiedlicher Hersteller zusammenarbeiten können
    \item \fatsf{Abstraktion:} Komplexe Vorgänge werden in verständliche Teile zerlegt
    \item \fatsf{Austauschbarkeit:} Einzelne Schichten können verändert werden, ohne das ganze System zu beeinflussen
\end{itemize}
\subsubsection*{Die 7 Schichten des OSI-Modells}
\begin{figure}[h]
    \centering
    \includegraphics[width=0.8\linewidth]{Chapter 1_Images/OSI Modell.jpeg}
\end{figure}
\begin{enumerate}
    \item \fatsf{Physical Layer (Physikalische Schicht)}\\
    \enquote{Wie werden einzelne Bits physisch übertragen?}
    \begin{itemize}
        \item Überträgt rohe Bits als elektrische Signale, Lichtimpulse oder Funkwellen
        \item Beispiel: Stecker, Kabel, Funkfrequenz, Spannung
        \item \fatsf{Hardware-orientiert}
        \item \fatsf{Protokoll-Beispiel:} RS232, Ethernet PHY
    \end{itemize}
    \item \fatsf{Data Link Layer (Sicherungsschicht)}\\
    \enquote{Wie machen wir aus den rohen Bits sinnvolle Datenrahmen und erkennen Fehler?}
    \begin{itemize}
        \item Erkennt \& behebt Übertragungsfehler (z.B. durch Prüfsummen)
        \item Teilt den Datenstrom in \fatsf{Frames}
        \item Flow-Control möglich (z.B. Sliding Window)
        \item Bei LANs:
        \begin{itemize}
            \item \fatsf{LLC (Logical Link Control):} Fehlererkennung
            \item \fatsf{MAC (Media Access Control):} Zugangskontrolle zum Übertragungsmedium (z.B. CSMA/CD)
        \end{itemize}
        \item \fatsf{Protokoll-Beispiel:} Ethernet, WLAN
    \end{itemize}
    \item \fatsf{Network Layer (Vermittlungsschicht)}\\
    \enquote{Wie kommt ein Paket von A nach B - auch über mehrere Netze hinweg?}
    \begin{itemize}
        \item Routing: Pfadwahl durch das Netz
        \item Adressierung (z.B. IP-Adressen)
        \item Fragmentierung \& Weiterleitung von Paketen
        \item \fatsf{Protokoll-Beispiel:} IP, X.25
    \end{itemize}
    \item \fatsf{Transport Layer (Transportschicht)}\\
    \enquote{Wie garantieren wir eine sichere Ende-zu-Ende-Kommunikation zwischen Prozessen?}
    \begin{itemize}
        \item Verbindungsaufbau \& -steuerung zwischen \fatsf{Anwendungen/Prozessen}
        \item Fehlerkontrolle, Wiederholungen, Reihenfolge
        \item Multiplexing von Datenstömen (mehrere Verbindungen über einen Kanal)
        \item \fatsf{Protokoll-Beispiel:} TCP (zuverlässig), UDP (schnell, aber unzuverlässig)
    \end{itemize}
        \begin{figure}[h]
        \centering
        \includegraphics[width=0.7\linewidth]{Chapter 1_Images/Transport Layer.jpeg}
    \end{figure}
    \item \fatsf{Session Layer (Sitzungssicht)}\\
    \enquote{Wie organisieren wir längere Kommunikationssitzungen?}
    \begin{itemize}
        \item Steuert Dialoge (z.B. wer darf wenn senden?)
        \item Checkpoints für Wiederherstellung nach Fehlern
        \item Heute selten genutzt, Funktionen oft in Transport-/Application-Schicht integriert
        \item \fatsf{Beispiel:} frühere Dienste wie Teletext
    \end{itemize}
    \item \fatsf{Presentation Layer (Darstellungsschicht)}\\
    \enquote{Wie sorgen wir dafür, dass beide Seiten die Daten verstehen?}
    \begin{itemize}
        \item Übersetzt Datenformate (z.B. Zeichencodierung: ASCII $\leftrightarrow$ Unicode)
        \item Syntaxvereinheitlichung
        \item \fatsf{Datenkompression \& Verschlüsselung} möglich
        \item \fatsf{Beispiel:} ASN.1, BER, google Protobuf
        \item In modernen Netzwerken oft Teil der Anwendungsschicht
    \end{itemize}
    \item \fatsf{Application Layer (Anwendungsschicht)}\\
    \enquote{Wie interagieren Anwendungen mit dem Netzwerk?}
    \begin{itemize}
        \item Bietet Netzwerkdienste für Benutzerprogramme
        \item Enthält viele \enquote{Service Elements} wie:
        \begin{itemize}
            \item Dateiübertragung (FTP)
            \item E-Mail (SMTP, X.400)
            \item Web (HTTP)
        \end{itemize}
        \item Praktisch: das, was du als Anwender siehst und nutzt
    \end{itemize} 
\end{enumerate}
\subsubsection*{Wichtige Hinweise}
\begin{itemize}
    \item Viele Dienste werden \fatsf{in mehreren Schichten} umgesetzt (z.B. Fehlerkorrektur)
    \item Das OSI-Modell ist \fatsf{konzeptionell} - in der Praxis dominiert das Internet-Modell
    \item Dennoch: \fatsf{OSI hilft enorm beim Verstehen komplexer Netzwerke}
\end{itemize}
\subsubsection*{Zwei Arten von Netzwerken / Protokollen}
Es gibt \fatsf{zwei grundlegende Arten}, wie Netzwerke Daten übertragen:
\begin{enumerate}
    \item \fatsf{Connection-Oriented (CO)}
    \begin{itemize}
        \item Verbindung wird vorher aufgebaut, wie bei einem Telefonat\\
        $\rightarrow$ Beispiel: \fatsf{TCP}
    \end{itemize} 
    \item \fatsf{Connectionless (CL)}
    \begin{itemize}
        \item Daten werden einfach losgeschickt, ohne vorherige Verbindung\\
        $\rightarrow$ Beispiel: \fatsf{UDP}
    \end{itemize}
\end{enumerate}
\subsubsection*{Conncection-Oriented (Verbindungsorientiert)}
\fatsf{Ablauf in 3 Phasen:}
\begin{enumerate}
    \item \fatsf{Verbindungsaufbau} (z.B. Handshake)
    \item \fatsf{Datenübertragung}
    \item \fatsf{Verbindungsabbau}
\end{enumerate}
Beispiel: \fatsf{Telefonanruf}
\begin{itemize}
    \item Erst wählt du die Nummer
    \item Dann sprichst du
    \item Dann legst du auf
\end{itemize}
\fatsf{Wichtig:}
\begin{itemize}
    \item Auch wenn der Nutzer \enquote{verbunden} ist, wissen die \fatsf{Router im Internet} davon \fatsf{nichts}
    \item Die Verbindung ist nur ein \fatsf{Zustand an beiden Enden} (also auf deinem Gerät \& dem Server)
    \item Das erlaubt bessere Dienste: \fatsf{Zuverlässigkeit, Flusskontrolle, Staukontrolle}
    \item \fatsf{Aber:} Das erzeugt \fatsf{Overhead} (mehr Daten, mehr Speicher, mehr Komplexität)
\end{itemize}
\subsubsection*{Connectionless (Verbindungsfrei)}
\fatsf{Ablauf:}
\begin{enumerate}
    \item \fatsf{Nur Datenübertragung} - keine Vorbereitung, keine Aufräumphase
\end{enumerate}
Beispiel: \fatsf{Postkarte verschicken}
\begin{itemize}
    \item Du schickst sie einfach ab, ohne zu fragen, ob jemand da ist
    \item Sie kommt vielleicht an - oder auch nicht
\end{itemize}
\fatsf{Eigenschaften:}
\begin{itemize}
    \item Kein Zustand nötig \fatsf{$\rightarrow$ Skalierbar!}
    \item Aber: \fatsf{Keine Garantie}, dass Daten ankommen oder geordnet sind
    \item Wird oft in Echtzeitanwendungen genutzt, wo Geschwindigkeit wichtig ist als Zuverlässigkeit (z.B. \fatsf{Video-Streaming, Online-Gaming})
\end{itemize}
\subsubsection*{Was ist besser?}
\begin{table}[H]
    \centering
    \rowcolors{2}{\thepagecolor}{fgcolor!10!\thepagecolor}
    \renewcommand{\arraystretch}{1.2}
    \begin{tabular}{lll}
        \toprule
        & \fatsf{Connection-Oriented (z.B. TCP)} & \fatsf{Connectionless (z.B. UDP)}\\
        \midrule
        Vorteil & Zuverlässig, ordentlich & Schnell, einfach, skalierbar\\
        Nachteil & Mehr Overhead & Keine Garantie\\
        Beispiel & Web, E-Mail & Video, DNS, Voice Chat\\
        \bottomrule 
    \end{tabular}
\end{table}
Du kannst sogar \fatsf{verbindungsorientierte Dienste auf verbindungslose Protokolle bauen} (z.B. QUIC auf UDP)
\subsubsection*{Beispiel: Was passiert, wenn du eine Webseite aufrufst}
Du öffnest den Browser $\rightarrow$ gibst eine URL ein $\rightarrow$ und dann passiert viel im Hintergrund
\begin{figure}[h]
    \centering
    \includegraphics[width=0.9\linewidth]{Chapter 1_Images/Layers in Action.jpeg}
\end{figure}
\begin{table}[H]
    \centering
    \rowcolors{2}{\thepagecolor}{fgcolor!10!\thepagecolor}
    \renewcommand{\arraystretch}{1.2}
    \begin{tabular}{ll}
        \toprule
        \fatsf{OSI-Schicht} & \fatsf{Aufgabe}\\
        \midrule
        \fatsf{7 (6,5) Anwendung} & HTTP-Anfrage erstellen, Antwort anzeigen\\
        \fatsf{4 - Transport} & TCP-Verbindung aufbauen, Daten in Reihenfolge bringen\\
        \fatsf{3 - Netzwerk} & IP: Routing durch Internet (Router zu Router)\\
        \fatsf{2 - Sicherung} & Ethernet/WLAN: Rahmen für jedes Paket\\
        \fatsf{1 - Bitübertragung} & Sendet die Bits über Kabel oder Funk\\
        \bottomrule 
    \end{tabular}
\end{table}